\documentclass[10pt,a4paper]{amsart}
\usepackage{amsmath,amssymb,amsthm,graphicx}
% Retained preamble for the Stacks Project chapter snapshot in this directory.
% The chapter's own preamble.tex is a copy of the shared Stacks Project preamble; the block
% below is sliced MECHANICALLY, line for line, from it. Nothing was added, removed or
% redefined; the wrapper class, the xr-hyper/hyperref wiring and the externaldocument list
% are omitted because the assembler supplies the document class.
%
% lines taken from preamble.tex: 3, 5, 9, 11, 12, 156, 157, 158, ... (51 lines total)
%
% package coverage restored from the chapter's own preamble (2026-10-13): xy, xy-2cell

\usepackage[all]{xy}\usepackage{verbatim}
\xyoption{2cell}
\UseAllTwocells

\newenvironment{reference}{\comment}{\endcomment}
\newenvironment{slogan}{\comment}{\endcomment}
\newenvironment{history}{\comment}{\endcomment}
\theoremstyle{plain}
\newtheorem{theorem}[subsection]{Theorem}
\newtheorem{proposition}[subsection]{Proposition}
\newtheorem{lemma}[subsection]{Lemma}
\theoremstyle{definition}
\newtheorem{definition}[subsection]{Definition}
\newtheorem{example}[subsection]{Example}
\newtheorem{exercise}[subsection]{Exercise}
\newtheorem{situation}[subsection]{Situation}
\theoremstyle{remark}
\newtheorem{remark}[subsection]{Remark}
\newtheorem{remarks}[subsection]{Remarks}
\numberwithin{equation}{subsection}
\def\lim{\mathop{\mathrm{lim}}\nolimits}
\def\colim{\mathop{\mathrm{colim}}\nolimits}
\def\Spec{\mathop{\mathrm{Spec}}}
\def\Hom{\mathop{\mathrm{Hom}}\nolimits}
\def\Ext{\mathop{\mathrm{Ext}}\nolimits}
\def\SheafHom{\mathop{\mathcal{H}\!\mathit{om}}\nolimits}
\def\SheafExt{\mathop{\mathcal{E}\!\mathit{xt}}\nolimits}
\def\Sch{\mathit{Sch}}
\def\Mor{\mathop{\mathrm{Mor}}\nolimits}
\def\Ob{\mathop{\mathrm{Ob}}\nolimits}
\def\Sh{\mathop{\mathit{Sh}}\nolimits}
\def\NL{\mathop{N\!L}\nolimits}
\def\CH{\mathop{\mathrm{CH}}\nolimits}
\def\proetale{{pro\text{-}\acute{e}tale}}
\def\etale{{\acute{e}tale}}
\def\QCoh{\mathit{QCoh}}
\def\Ker{\mathop{\mathrm{Ker}}}
\def\Im{\mathop{\mathrm{Im}}}
\def\Coker{\mathop{\mathrm{Coker}}}
\def\Coim{\mathop{\mathrm{Coim}}}
\DeclareMathSymbol{\boxtimes}{\mathbin}{AMSa}{"02}
\def\QCohstack{\mathcal{QC}\!\mathit{oh}}
\def\Cohstack{\mathcal{C}\!\mathit{oh}}
\def\Spacesstack{\mathcal{S}\!\mathit{paces}}
\def\Quotfunctor{\mathrm{Quot}}
\def\Hilbfunctor{\mathrm{Hilb}}
\def\Curvesstack{\mathcal{C}\!\mathit{urves}}
\def\Polarizedstack{\mathcal{P}\!\mathit{olarized}}
\def\Complexesstack{\mathcal{C}\!\mathit{omplexes}}
\def\Pic{\mathop{\mathrm{Pic}}\nolimits}
\def\Picardstack{\mathcal{P}\!\mathit{ic}}
\def\Picardfunctor{\mathrm{Pic}}
\def\Deformationcategory{\mathcal{D}\!\mathit{ef}}
\usepackage{cleveref}
\crefname{originallemma}{original lemma}{original lemmas}
\Crefname{originallemma}{Original lemma}{Original lemmas}
\makeatletter
\newlabel{lemma-affine-morphism}{{08II}{source}{}{sourcepointer-1}{}}
\newlabel{lemma-affine-morphism@cref}{{[originallemma][0][]08II}{[1][0][]source}{}{}{}}
\makeatother
\crefname{originallemma}{original lemma}{original lemmas}
\Crefname{originallemma}{Original lemma}{Original lemmas}
\makeatletter
\newlabel{lemma-affine-morphism-pull-push}{{08IJ}{source}{}{sourcepointer-2}{}}
\newlabel{lemma-affine-morphism-pull-push@cref}{{[originallemma][0][]08IJ}{[1][0][]source}{}{}{}}
\makeatother
\crefname{originallemma}{original lemma}{original lemmas}
\Crefname{originallemma}{Original lemma}{Original lemmas}
\makeatletter
\newlabel{lemma-cohomology-base-change}{{08IN}{source}{}{sourcepointer-3}{}}
\newlabel{lemma-cohomology-base-change@cref}{{[originallemma][0][]08IN}{[1][0][]source}{}{}{}}
\makeatother
\crefname{originallemma}{original lemma}{original lemmas}
\Crefname{originallemma}{Original lemma}{Original lemmas}
\makeatletter
\newlabel{lemma-quasi-coherence-direct-image}{{08FA}{source}{}{sourcepointer-4}{}}
\newlabel{lemma-quasi-coherence-direct-image@cref}{{[originallemma][0][]08FA}{[1][0][]source}{}{}{}}
\makeatother
\crefname{originallemma}{original lemma}{original lemmas}
\Crefname{originallemma}{Original lemma}{Original lemmas}
\makeatletter
\newlabel{lemma-quasi-coherence-pullback}{{08F4}{source}{}{sourcepointer-5}{}}
\newlabel{lemma-quasi-coherence-pullback@cref}{{[originallemma][0][]08F4}{[1][0][]source}{}{}{}}
\makeatother
\crefname{originallemma}{original lemma}{original lemmas}
\Crefname{originallemma}{Original lemma}{Original lemmas}
\makeatletter
\newlabel{lemma-quasi-coherence-tensor-product}{{08F5}{source}{}{sourcepointer-6}{}}
\newlabel{lemma-quasi-coherence-tensor-product@cref}{{[originallemma][0][]08F5}{[1][0][]source}{}{}{}}
\makeatother
\crefname{originallemma}{original lemma}{original lemmas}
\Crefname{originallemma}{Original lemma}{Original lemmas}
\makeatletter
\newlabel{lemma-tor-independent}{{08IQ}{source}{}{sourcepointer-7}{}}
\newlabel{lemma-tor-independent@cref}{{[originallemma][0][]08IQ}{[1][0][]source}{}{}{}}
\makeatother
\crefname{originallemma}{original lemma}{original lemmas}
\Crefname{originallemma}{Original lemma}{Original lemmas}
\makeatletter
\newlabel{more-algebra-lemma-tor-independent}{{08HX}{source}{}{sourcepointer-8}{}}
\newlabel{more-algebra-lemma-tor-independent@cref}{{[originallemma][0][]08HX}{[1][0][]source}{}{}{}}
\makeatother
\crefname{originallemma}{original lemma}{original lemmas}
\Crefname{originallemma}{Original lemma}{Original lemmas}
\makeatletter
\newlabel{sites-cohomology-remark-base-change}{{07A7}{source}{}{sourcepointer-9}{}}
\newlabel{sites-cohomology-remark-base-change@cref}{{[originallemma][0][]07A7}{[1][0][]source}{}{}{}}
\makeatother
\crefname{originallemma}{original lemma}{original lemmas}
\Crefname{originallemma}{Original lemma}{Original lemmas}
\makeatletter
\newlabel{spaces-morphisms-lemma-base-change-affine}{{03WI}{source}{}{sourcepointer-10}{}}
\newlabel{spaces-morphisms-lemma-base-change-affine@cref}{{[originallemma][0][]03WI}{[1][0][]source}{}{}{}}
\makeatother
\title{Stacks, Derived Categories of Spaces: compare-base-change}
\author{The Stacks Project Authors}
\date{stacks-project:spaces-perfect.tex:local-snapshot-20261012; GNU FDL 1.2 or later}
\begin{document}
\maketitle
\noindent\textit{Editorial note.} Statement span is lines 5079-5095 of sources/stacks-spaces-perfect/pinned-original-source.tex; proof span is lines 5097-5135, the complete author proof (37 body lines).\par\medskip
\noindent\textit{Editorial note.} Navigation stubs: 10 cited-but-not-included labels are supplied as non-load-bearing pointers with their pinned tags where the snapshot tag index has them.\par\medskip
\noindent\textit{Editorial note.} extra\_packages is empty on purpose: the retained preamble already loads xy with options and a second declaration clashes in the assembler wrapper.\par\medskip
\noindent\textit{Editorial note.} Text is retained verbatim, including the author's own typos and placeholders.\par\medskip
\noindent\textit{Editorial note.} Queue index 31 of queue\_spaceperfect.json (pb=1846, stmt 5079-5095, proof 5097-5135).\par\medskip
\noindent\textit{Editorial note.} stacks-spaces-perfect wave sp01 (2026-10-14): fresh adjudication accept per lane-s3/sp01\_judgments; pool per Chief 12:4x unfreeze (reserve-pool workflow, six-check clean); spaces-perfect.tex 126 with-proof rows, 70 b107-occupied labels avoided, 56 free rows judged 32 accept / 24 reject.\par\medskip
\noindent\textit{Editorial note.} Cross-chapter same-label warning (NOT an occupancy conflict, Chief ruling 2026-10-13 evening): label 'lemma-compare-base-change' also exists in stacks-perfect — a different proposition in a different chapter file.\par\medskip
\noindent\textit{Allowed standard prerequisites (editorial).}\begin{itemize}
\item Explicit prerequisite boundary. The cited inputs are prerequisites; the argument in the retained span is not replaced by any cited result.
\end{itemize}\medskip
\ifdefined\hypertarget\hypertarget{sourcepointer-1}{}\fi
\noindent\textit{Source reference.} Original lemma 08II: sources/stacks-spaces-perfect/pinned-original-source.tex, source lines 965--970. Cited by the retained statement or proof but not part of the collected spans; supplied only as a navigation pointer. The argument does not depend on its content.\par\medskip
\ifdefined\hypertarget\hypertarget{sourcepointer-2}{}\fi
\noindent\textit{Source reference.} Original lemma 08IJ: sources/stacks-spaces-perfect/pinned-original-source.tex, source lines 985--989. Cited by the retained statement or proof but not part of the collected spans; supplied only as a navigation pointer. The argument does not depend on its content.\par\medskip
\ifdefined\hypertarget\hypertarget{sourcepointer-3}{}\fi
\noindent\textit{Source reference.} Original lemma 08IN: sources/stacks-spaces-perfect/pinned-original-source.tex, source lines 4883--4892. Cited by the retained statement or proof but not part of the collected spans; supplied only as a navigation pointer. The argument does not depend on its content.\par\medskip
\ifdefined\hypertarget\hypertarget{sourcepointer-4}{}\fi
\noindent\textit{Source reference.} Original lemma 08FA: sources/stacks-spaces-perfect/pinned-original-source.tex, source lines 790--804. Cited by the retained statement or proof but not part of the collected spans; supplied only as a navigation pointer. The argument does not depend on its content.\par\medskip
\ifdefined\hypertarget\hypertarget{sourcepointer-5}{}\fi
\noindent\textit{Source reference.} Original lemma 08F4: sources/stacks-spaces-perfect/pinned-original-source.tex, source lines 620--625. Cited by the retained statement or proof but not part of the collected spans; supplied only as a navigation pointer. The argument does not depend on its content.\par\medskip
\ifdefined\hypertarget\hypertarget{sourcepointer-6}{}\fi
\noindent\textit{Source reference.} Original lemma 08F5: sources/stacks-spaces-perfect/pinned-original-source.tex, source lines 645--650. Cited by the retained statement or proof but not part of the collected spans; supplied only as a navigation pointer. The argument does not depend on its content.\par\medskip
\ifdefined\hypertarget\hypertarget{sourcepointer-7}{}\fi
\noindent\textit{Source reference.} Original lemma 08IQ: sources/stacks-spaces-perfect/pinned-original-source.tex, source lines 4965--4986. Cited by the retained statement or proof but not part of the collected spans; supplied only as a navigation pointer. The argument does not depend on its content.\par\medskip
\ifdefined\hypertarget\hypertarget{sourcepointer-8}{}\fi
\noindent\textit{Source reference.} Original lemma 08HX: sources/stacks-spaces-perfect/pinned-original-source.tex, source lines 5135--5135. Cited by the retained statement or proof but not part of the collected spans; supplied only as a navigation pointer. The argument does not depend on its content.\par\medskip
\ifdefined\hypertarget\hypertarget{sourcepointer-9}{}\fi
\noindent\textit{Source reference.} Original lemma 07A7: sources/stacks-spaces-perfect/pinned-original-source.tex, source lines 5135--5135. Cited by the retained statement or proof but not part of the collected spans; supplied only as a navigation pointer. The argument does not depend on its content.\par\medskip
\ifdefined\hypertarget\hypertarget{sourcepointer-10}{}\fi
\noindent\textit{Source reference.} Original lemma 03WI: sources/stacks-spaces-perfect/pinned-original-source.tex, source lines 5135--5135. Cited by the retained statement or proof but not part of the collected spans; supplied only as a navigation pointer. The argument does not depend on its content.\par\medskip
\begin{lemma}
\label{lemma-compare-base-change}
Let $S$ be a scheme. Let $g : Y' \to Y$ be a morphism of algebraic spaces over
$S$. Let $f : X \to Y$ be a quasi-compact and quasi-separated morphism of
algebraic spaces over $S$. Consider the base change diagram
$$
\xymatrix{
X' \ar[r]_{g'} \ar[d]_{f'} &
X \ar[d]^f \\
Y' \ar[r]^g &
Y
}
$$
If $X$ and $Y'$ are Tor independent over $Y$, then for all
$E \in D_\QCoh(\mathcal{O}_X)$ we have
$Rf'_*L(g')^*E = Lg^*Rf_*E$.
\end{lemma}

\begin{proof}
For any object $E$ of $D(\mathcal{O}_X)$ we can use
Cohomology on Sites, Remark \ref{sites-cohomology-remark-base-change}
to get a canonical base change map $Lg^*Rf_*E \to Rf'_*L(g')^*E$. To check this
is an isomorphism we may work \'etale locally on $Y'$. Hence we may assume
$g : Y' \to Y$ is a morphism of affine schemes. In particular, $g$
is affine and it suffices to show that
$$
Rg_*Lg^*Rf_*E \to Rg_*Rf'_*L(g')^*E = Rf_*(Rg'_* L(g')^* E)
$$
is an isomorphism, see Lemma \ref{lemma-affine-morphism}
(and use Lemmas \ref{lemma-quasi-coherence-pullback},
\ref{lemma-quasi-coherence-tensor-product}, and
\ref{lemma-quasi-coherence-direct-image}
to see that the objects $Rf'_*L(g')^*E$ and $Lg^*Rf_*E$
have quasi-coherent cohomology sheaves). Note that $g'$ is
affine as well (Morphisms of Spaces, Lemma
\ref{spaces-morphisms-lemma-base-change-affine}).
By Lemma \ref{lemma-affine-morphism-pull-push} the map becomes a map
$$
Rf_*E \otimes_{\mathcal{O}_Y}^\mathbf{L} g_*\mathcal{O}_{Y'}
\longrightarrow
Rf_*(E \otimes_{\mathcal{O}_X}^\mathbf{L} g'_*\mathcal{O}_{X'})
$$
Observe that $g'_*\mathcal{O}_{X'} = f^*g_*\mathcal{O}_{Y'}$. Thus by
Lemma \ref{lemma-cohomology-base-change} it suffices to prove that
$Lf^*g_*\mathcal{O}_{Y'} = f^*g_*\mathcal{O}_{Y'}$. This follows from our
assumption that $X$ and $Y'$ are Tor independent over $Y$. Namely, to
check it we may work \'etale locally on $X$, hence we may also assume $X$
is affine. Say $X = \Spec(A)$, $Y = \Spec(R)$ and $Y' = \Spec(R')$.
Our assumption implies that $A$ and $R'$ are Tor independent over $R$
(see
Lemma \ref{lemma-tor-independent}
and
More on Algebra, Lemma \ref{more-algebra-lemma-tor-independent}), i.e.,
$\text{Tor}_i^R(A, R') = 0$ for $i > 0$. In other words
$A \otimes_R^\mathbf{L} R' = A \otimes_R R'$ which exactly means
that $Lf^*g_*\mathcal{O}_{Y'} = f^*g_*\mathcal{O}_{Y'}$.
\end{proof}

\end{document}
